\subsection{Requisiti qualitativi}
Software con le stesse funzionalità ma qualità differenti sono considerati software differenti. Per tanto i requisiti qualitativi sono uno dei fattori che determinano la qualità generale del software. \\
I requisiti qualitativi sono classificati in:
\begin{itemize}
    \item \textbf{di prodotto}, relativi all'erogazione del servizio
    \item \textbf{organizzativi}, per lo sviluppo interno all'azienda
    \item \textbf{esterni}, basati sul contesto generale dell'implementazione
\end{itemize}
I requisiti qualitativi si basano su \textbf{metriche di valutazione}, cioè funzioni che forniscono un'unità di misura. Alcuni esempi di qualità sono elencati nell'\hyperref[appendix:qualities]{approndimento teorico}.
